\chapter{Metodologia de pesquisa}\label{cap:metodologia}

Essa seção existirá caso o estudante e seu orientador entendam que o conjunto de métodos de pesquisa utilizados no trabalho devem ser detalhados e relatados a fim de que o leitor acompanhe como a investigação científica foi desenvolvida. Essa seção apresenta os métodos de investigação utilizados, a ordem em que aparecem no texto e as razões para escolha de cada um deles, ou seja, qual o objetivo esperado com cada método.

Uma pesquisa pode fazer uso de um único método. E em alguns casos, orientador e estudante decidem colocar essa informação referente aos métodos de investigação em uma subseção da introdução.

A seção de metodologia difere da seção de resultados. A primeira aborda como os dados serão coletados e analisados, com qual método ou abordagem. A segunda relata os resultados advindos dessa coleta e dessa análise. Por exemplo, uma seção de análise demográfica dos dados deverá estar entre os resultados. 

São vários os métodos de pesquisa: estudo de caso, experimento, levantamento estatístico, fenomenologia, hermenêutica, etnografia, design science research, pesquisa-ação, ciência da implementação etc. O estudante deverá escolher os métodos que melhor se adequem ao tipo de contribuição que quer oferecer e que dialoguem com as pesquisas anteriores sobre o tema, permitindo que ele faça comparações com trabalhos correlatos. 

A Metodologia de Pesquisa pode ser dividida em subseções para melhor compreensão do leitor. Dependendo de cada método escolhido, poderá haver diferentes subseções. Segue uma sugestão que deverá se adequar ao estilo do estudante e de seu orientador.

\section{Tipo de pesquisa}
Quanto aos objetivos: sua pesquisa é exploratória, descritiva ou explicativa? Por quê?
Quanto aos meios: bibliográfica, documental, estudo de caso, etc... 

\section{Seleção dos sujeitos / Entrevistados}
Quantos são os entrevistados? Qual o cargo que ocupam e há quanto tempo? Por que é importante entrevistá-los?
Caso seja uma pesquisa mais objetiva: qual a população e a amostra? Qual a representatividade esperada?

\section{Coleta de dados}
Como será feita a coleta dos dados? A partir de documentação direta (observação do fenômeno ou entrevistas) ou indireta (leitura de dados sobre o fenômeno)? Se houver roteiro ou questionário, adicionar ao anexo.

\section{Análise de dados}
Como o pesquisador pretende fazer as análises dos dados que coletará? Usará alguma abordagem interpretativa caso a pesquisa seja qualitativa? Usará quais ferramentas de análise estatística caso esteja atuando com métodos quantitativos.